\section{Related Work}
\label{sec:related}
\para{Leakage and suppression.}
\citet{holtzman2026secret} study involuntary information leakage in writing, including secret/no-secret detection and discrimination between different secrets. Their reported absence of literal secret mentions makes our frequent exact disclosures a material protocol difference. We adapt their detection task, word inventory, and decoy precedent; we do not claim a clean replication of their nonliteral effect. Semantic leakage also occurs when irrelevant prompt concepts influence outputs without an explicit suppression instruction \cite{gonen2024semantic}. Ironic-negation experiments examine how load changes suppression behavior \cite{mann2025whitebear}, while \citet{ramnauth2026attentional} find that suppressed concepts remain decodable. These studies motivate separating literal compliance, semantic influence, and residual information.

\para{Planning stories and future tokens.}
Plan-and-Write separates storylines from realization \cite{yao2018plan}; Re3 and DOC use recursive drafting and detailed outlines to improve long-story generation \cite{yang2022re3,yang2022doc}. Iterative planning also supports suspenseful writing \cite{xie2024suspense}. These objectives do not imply secrecy: a coherent or suspenseful story may deliberately foreshadow its ending. At the representation level, response attributes can be predicted before generation \cite{dong2025planning}. \citet{wu2024future} distinguish anticipatory computation from representations that incidentally help future prediction. Our outline/context comparison addresses behavioral disclosure, and our probes do not identify anticipatory computation specifically.

\para{Secret readouts and controlled interventions.}
Taboo model organisms and broader secret-knowledge benchmarks compare black-box elicitation with white-box readouts \cite{cywinski2025latent,cywinski2025eliciting}. Their tasks include learned secrets and deliberate hint generation, whereas our secret stays in the prompt during unrelated fiction. Sparse activation editing studies instruction compliance in narratives \cite{zhao2025editing}; contextual-privacy steering studies information-flow norms \cite{wang2026privacy}. These motivate target, random, and other-concept controls together with utility checks. We retain a failed gate rather than infer an intervention effect from offline decoding.

\para{Confidentiality and private state.}
Contextual integrity benchmarks test whether models respect appropriate information flows across social settings \cite{mireshghallah2023privacy}. Private-memory work studies consistency when agents must preserve hidden choices across interaction histories \cite{baldelli2026hangman}. Neither task is equivalent to suppressing an arbitrary word while its original prompt remains available. Our results concern this limited writing protocol and do not certify practical confidentiality or private memory.
